% =====================================================================
\section{Kernel derivation}\label{sec:derivation}
% =====================================================================

This section derives $\kappa(q^{2})$ in three layers: the continuum
$\varphi$-tuned closure operator (\ref{subsec:layer1}), the bounded
Dirichlet eigenmode (\ref{subsec:layer2}), and the discrete
2I-equivariant graph realisation on the 600-cell $V_{600}$
(\ref{subsec:layer3}). The fit form
$\Delta C_{9}^{\mathrm{VFD}}(q^{2})=-A\kappa(q^{2})$ uses the Layer-3
kernel; the continuum and bounded variants are reported as theoretical
controls in section~\ref{sec:results}.

\subsection{Layer 1 — continuum $\varphi$-tuned Green's function}\label{subsec:layer1}

Consider the one-dimensional closure operator
\begin{equation}\label{eq:Lphi}
L_{\varphi} \;=\; -\frac{\partial^{2}}{\partial x^{2}} + \frac{1}{\varphi^{2}},
\end{equation}
on the dimensionless coordinate
$x = (q^{2}-q^{2}_{\mathrm{mid}})/\Delta_{\psi}$, with
$q^{2}_{\mathrm{mid}}=\tfrac{1}{2}(m^{2}_{J/\psi}+m^{2}_{\psi(2S)})
\approx 11.59\,\mathrm{GeV}^{2}$ and
$\Delta_{\psi}=\tfrac{1}{2}(m^{2}_{\psi(2S)}-m^{2}_{J/\psi})\approx
4.00\,\mathrm{GeV}^{2}$. The free-space Green's function $G(x;0)$
satisfying $L_{\varphi}\, G(x;0)=\delta(x)$ is
\begin{equation}\label{eq:greens_continuum}
G(x;0) = \frac{\varphi}{2}\, e^{-|x|/\varphi},
\end{equation}
giving the empirical kernel
$\kappa_{\mathrm{cont}}(x) = e^{-|x|/\varphi}$ which served as the
champion shape in the project's initial single-parameter scan over
candidate $q^{2}$ kernels. The choice of mass parameter
$\varphi^{-1}$ in Eq.~\eqref{eq:Lphi} is the only continuous
constant in this paper that is not data-fit. It is the unique scale
that ties this operator to the binary-icosahedral-symmetric graph
$V_{600}$ used in the discrete realisation below (the vertex
coordinates of $V_{600}$ are built from $\varphi$ as the only
irrational unit; \S\ref{subsec:layer3}); the wider framework context
that motivated this choice is described in companion preprints
\citep{SmartCrystallisation2026, SmartAdaptiveClosureTransport2026},
but is not relied on for any quantitative claim in this paper.

\subsection{Layer 2 — bounded Dirichlet eigenmode}\label{subsec:layer2}

Restricting $L_{\varphi}$ to the kinematic interval $[-x_{\max},\,
+x_{\max}]$ with Dirichlet conditions $\psi(\pm x_{\max})=0$ and
$x_{\max}=2.886$ (the rescaled
$q^{2}\in[0,19]\,\mathrm{GeV}^{2}$ window) gives the spectral problem
\begin{equation}
L_{\varphi}\, \psi_{n}(x) = \lambda_{n}\, \psi_{n}(x),
\qquad
\psi_{n}(\pm x_{\max}) = 0,
\end{equation}
with the lowest eigenmode
\begin{equation}\label{eq:dirichlet_mode}
\psi_{1}(x) = \cos\!\Bigl(\frac{\pi x}{2\,x_{\max}}\Bigr),
\qquad
\lambda_{1} = \frac{1}{\varphi^{2}} + \Bigl(\frac{\pi}{2 x_{\max}}\Bigr)^{2}.
\end{equation}
Both single-mode shapes (\ref{eq:greens_continuum}) and
(\ref{eq:dirichlet_mode}) fit the LHCb 2025 $P_{5}'$ data at
$\Delta\mathrm{AIC}\approx -7$ vs $\mathrm{FREE\_C9}$ on $P_{5}'$ alone
under the linearised \texttt{flavio} response (section~\ref{sec:results}).
A 2-mode Dirichlet truncation with two amplitudes
$(c_{1}, c_{3})$ is not preferred over the single mode after the
AIC penalty: $c_{3}/c_{1} = 0.18$, well below any $\varphi$-rational
threshold. The data is rank-1 in this basis on $P_{5}'$.

\subsection{Layer 3 — discrete 2I-equivariant lift on $V_{600}$}\label{subsec:layer3}

The continuum kernel of Eq.~\eqref{eq:greens_continuum} is recovered as
the shell-mean projection of a discrete Green's response on a finite
graph carrying the binary icosahedral symmetry $2I$. We build this
graph and verify its symmetry intrinsically.

\paragraph{Vertex set.} The 120 vertices of the regular 600-cell
$V_{600}\subset\mathbb{R}^{4}$ form the binary icosahedral group $2I$
under unit-quaternion multiplication
\citep{Coxeter1973RegularPolytopes,Conway2003Quaternions}. They are
generated by the eight unit basis vectors $\pm e_{i}$, the sixteen
half-integer vectors $\tfrac{1}{2}(\pm 1, \pm 1, \pm 1, \pm 1)$, and
the ninety-six even-permutation vectors of
$\tfrac{1}{2}(0, \pm 1, \pm \varphi^{-1}, \pm \varphi)$, totalling
$8+16+96 = 120$. They occupy the unit 3-sphere $S^{3}$.

\paragraph{Edge set.} Two vertices $v, w\in V_{600}$ are adjacent in
the 1-skeleton iff $\langle v, w\rangle = \varphi/2$. This produces a
12-regular graph with exactly 720 edges, in agreement with the standard
600-cell construction.

\paragraph{Inner-product shells.} Fixing a base vertex $v_{0}$ (we use
$v_{0}=e_{1}$), the function $\langle v, v_{0}\rangle$ takes nine
distinct values across $V_{600}$:
\begin{equation}\label{eq:shells}
\langle v, v_{0}\rangle \in
\{+1, +\tfrac{\varphi}{2}, +\tfrac{1}{2}, +\tfrac{1}{2\varphi}, 0,
   -\tfrac{1}{2\varphi}, -\tfrac{1}{2}, -\tfrac{\varphi}{2}, -1\}
\end{equation}
with shell sizes
$\{1, 12, 20, 12, 30, 12, 20, 12, 1\}$, summing to 120 and matching the
inner-product spectrum derived from the binary icosahedral group's
conjugacy classes. The shell index $s\in\{0,\dots,8\}$ provides the
discrete proxy for the dimensionless coordinate $x$ via
\begin{equation}\label{eq:x_shell}
x_{s} \;=\; (s - s_{\mathrm{mid}})\, \frac{x_{\max}}{s_{\mathrm{mid}}},
\qquad s_{\mathrm{mid}}=4.
\end{equation}

\paragraph{Discrete Green's response.} The standard graph Laplacian
$L_{V_{600}} = D - A$ with $D$ the degree matrix and $A$ the adjacency
matrix, regularised by the framework's discrete mass $\varphi^{-2}$,
gives the response
\begin{equation}\label{eq:psi_discrete}
\psi \;=\; \bigl(L_{V_{600}} + \varphi^{-2} I\bigr)^{-1}\, f,
\qquad
f_{v} \;=\; \begin{cases}1/|S_{4}| & v\in S_{4} \\ 0 & \text{otherwise}\end{cases}
\end{equation}
with $S_{4}$ the equatorial inner-product shell (30 vertices).

\paragraph{Shell-mean projection.} Averaging $\psi$ over each shell
gives the nine-component shell-mean profile
$\overline{\psi}_{s} = |S_{s}|^{-1}\sum_{v\in S_{s}}\psi_{v}$.
Linear interpolation in $x_{s}$ gives the continuous kernel
\begin{equation}\label{eq:kappa_lift}
\kappa_{\mathrm{V_{600}}}(x) \;=\; \mathrm{interp}\!\bigl(\{x_{s}\},
   \{\overline{\psi}_{s}\}\bigr)(x).
\end{equation}
The empirical correlation between Eqs.~\eqref{eq:greens_continuum} and
\eqref{eq:kappa_lift} on the LHCb 2025 bin centres is $r=0.997$. The
discrete kernel is the object that enters Eq.~\eqref{eq:fit_form} in
all subsequent fits; it is plotted in Figure~\ref{fig:F1_kernel_shape}
together with the continuum exponential and the LHCb 2025 bin
centres.

\begin{figure}[h]
\centering
\includegraphics[width=0.92\textwidth]{fig_F1_kernel_shape.pdf}
\caption{The geometry-derived response kernel $\kappa(q^{2})$
plotted over the LHCb $q^{2}$ window. Solid blue: the discrete
$V_{600}$ shell-mean of Eq.~\eqref{eq:kappa_lift} (Layer 3, used in
all fits). Dashed grey: the continuum Green's function
$e^{-|x|/\varphi}$ from Layer 1, with empirical correlation $r=0.997$
to the discrete shell-mean on the LHCb bin centres (red points).
Vertical dotted lines mark the $J/\psi$ and $\psi(2S)$ resonance
windows, which are excluded from the fits.
Reproducible from \texttt{scripts/wo017\_paper\_figures.py}.}
\label{fig:F1_kernel_shape}
\end{figure}

\paragraph{Cocycle, edge weighting, and variant selection.} The
framework's pentagonal cocycle
$\kappa(v) = (s(v)-s_{\mathrm{mid}})^{2}\in\{0,1,4,9,16\}$ provides a
vertex potential $\omega_{+}=\varphi^{\kappa(v)}$. Three discrete
edge-weighting schemes are admissible: unweighted ($w_{vw}=1$),
$\varphi$-cocycle geometric mean
($w_{vw}=\sqrt{\omega_{+}(v)\,\omega_{+}(w)}$), and $\varphi$-cocycle
arithmetic mean ($w_{vw}=\tfrac{1}{2}[\omega_{+}(v)+\omega_{+}(w)]$).
The discrete Green's response Eq.~\eqref{eq:psi_discrete} is computed
for each variant; the shell-mean profile is then compared to the
continuum kernel $e^{-|x|/\varphi}$ and to the LHCb 2025 $P_{5}'$ data
on the LHCb bin grid. Two criteria are used: a \emph{pure-geometry}
criterion (correlation between the discrete shell-mean and the Layer-1
continuum kernel, no LHCb data involvement), and a \emph{data}
criterion ($\chi^{2}$ against the LHCb 2025 single-observable fit).

\begin{table}[h]
\centering
\small
\begin{tabular}{lccc}
\toprule
variant & $\mathrm{corr}(\overline{\psi}, e^{-|x|/\varphi})$ &
        $\chi^{2}$ vs LHCb 2025 $P_{5}'$ &
        ranking \\
\midrule
unweighted Laplacian          & $\mathbf{0.9968}$ & $\mathbf{13.555}$ & $1$ on both \\
$\varphi$-cocycle geometric    & $0.9130$ & $14.713$ & $2$ on both \\
$\varphi$-cocycle arithmetic   & $0.8989$ & $14.782$ & $3$ on both \\
\bottomrule
\end{tabular}
\caption{Variant selection. The unweighted Laplacian wins on both the
pure-geometry criterion (correlation with the Layer-1 continuum
kernel) and the LHCb-data criterion. The two rankings agree, so the
variant choice is consistent with selection on geometric grounds
independent of the data; the LHCb data only confirms it.
Reproducible from \texttt{scripts/wo016b\_variant\_geometry.py}
(reports/wo016b\_variant\_geometry.csv).}
\label{tab:variant_selection}
\end{table}

The unweighted Laplacian is selected as the primary kernel. The
two-criterion agreement means the variant choice does not consume an
extra effective fit parameter: the variant choice would have been the
same if only the geometric criterion had been applied, and the AIC
counting $k=1$ for VFD\_GREEN\_600CELL is therefore retained
throughout this paper. The historical sequence is acknowledged in
\S\ref{sec:limitations}: the data was looked at first, and only later
verified to agree with the pure-geometry ranking.

\paragraph{Spectral content.} An eigenvalue decomposition of
$L_{V_{600}}$ produces 120 modes; truncating $\psi$ to the lowest
$N$ eigenmodes and reconstructing $\kappa$ shows that the data-relevant
information lives in a single isotypic block at $\lambda=9.0$ with
multiplicity 6. Truncations below $N=6$ degrade the fit; truncations
at $N\geq 6$ recover the full result. This rank-6 structural
compression is reported as a spectral diagnostic
(\texttt{src/vfd\_b\_anomaly/wo011\_spectral.py} in the project
repository); the cross-dataset fit itself uses the full Green's
response.

\paragraph{Self-contained reading.} The construction in this section
is self-contained: it requires only the standard graph Laplacian
(\S\ref{subsec:layer3}), the equatorial-shell source
(Eq.~\ref{eq:psi_discrete}), the $\varphi^{-2}$ regulariser
(Eq.~\ref{eq:Lphi}), and the variant table
(Table~\ref{tab:variant_selection}). A reader unfamiliar with the
broader VFD framework can take the kernel as defined here and
proceed to \S\ref{sec:results}. Two structural extensions — the
proper Hodge edge Laplacian $L_{\mathrm{edge}} =
\delta_{2}\delta_{2}^{\mathsf T} + \delta_{1}^{\mathsf T}\delta_{1}$
with 2I-equivariant boundary maps from the 1200 triangular faces,
and the Lyapunov-flow selection rule discussed in
\citep{SmartAdaptiveClosureTransport2026} — are noted in
\S\ref{sec:limitations} as deferred but are not used here.
